\documentclass[a4paper,twoside]{article}

\usepackage{morewrites}

\usepackage[svgnames,dvipsnames,x11names]{xcolor}
\usepackage{graphicx}
\usepackage[english]{babel}
\usepackage[colorlinks=true, linkcolor=NavyBlue, urlcolor=NavyBlue, citecolor=NavyBlue]{hyperref}
\usepackage[a4paper]{geometry}
\usepackage{ts-fonts}
\usepackage{booktabs}
\usepackage{tabularx}
\usepackage{amsmath}
\usepackage{float}
\usepackage{seqsplit}
\newcolumntype{Y}{>{\centering\arraybackslash}X}
\newcolumntype{Z}{>{\raggedleft\arraybackslash}X}
\usepackage{ts-code}

\usepackage{lastpage}
\usepackage{refcount}
\makeatletter
\def\ps@plain{
  \let\@oddhead\@empty
  \let\@evenhead\@empty
  \def\@oddfoot{
    \hfil
    \ifnum\getpagerefnumber{LastPage}>1\relax
      \thepage
    \fi
    \hfil}
  \let\@evenfoot\@oddfoot
}
\makeatother

\title{TeXSmith Extensions}
\author{}\date{}

\begin{document}

\maketitle

\thispagestyle{plain}
TeXSmith ships its Markdown extensions directly inside the \texttt{texsmith}
distribution. Once \texttt{pip install texsmith} is done, you can use them via the
extension modules (\texttt{texsmith.extensions.smallcaps}, \texttt{texsmith.extensions.index},
...) or the entry-point aliases registered in \texttt{pyproject.toml}. The exact set
enabled by the conversion pipeline is \texttt{texsmith.adapters.markdown.DEFAULT\_MARKDOWN\_EXTENSIONS}.

Every extension follows the same structure:

\begin{itemize}
\item{} A Python-Markdown class or \texttt{makeExtension()} factory available as
  \texttt{texsmith.extensions.<name>}.
\item{} Optional reader lowerings (\texttt{@reads}) and writer emitters (\texttt{@writes}) that
  teach the IR pipeline how to deal with the extra HTML nodes created by the
  Markdown layer (see \hyperref[snippet:<HASH>]{Readers \& Writers}).
\item{} Optional MkDocs plugins that keep search indexes in sync.

\end{itemize}
\section{Built-in extensions}\label{built-in-extensions}

In order to align Markdown with \LaTeX{} capabilities, TeXSmith provides the
following extensions under the \texttt{texsmith} namespace:

\begin{center}
\begin{tabularx}{\linewidth}{>{\raggedright\arraybackslash}X>{\raggedright\arraybackslash}X}
\toprule
\textbf{Module} & \textbf{Purpose} \\

\midrule
\texttt{texsmith.extensions.smallcaps} & \texttt{\_\_text\_\_} syntax mapped to \texttt{<span class="texsmith-\allowbreak{}smallcaps">}. \\
\texttt{texsmith.extensions.latex\_raw} & \texttt{/// latex} fences and \texttt{\{latex\}[x]} inline snippets injected as hidden HTML. \\
\texttt{texsmith.extensions.latex\_text} & Styles the literal \texttt{LaTeX} token in running text. \\
\texttt{texsmith.extensions.missing\_footnotes} & Warns about references to undefined footnotes. \\
\texttt{texsmith.extensions.multi\_citations} & Normalises \texttt{[\^{}foo,bar]} blocks to footnotes. \\
\texttt{texsmith.extensions.mermaid} & Inlines Mermaid diagrams pointed to by Markdown images. \\
\texttt{texsmith.extensions.texlogos} & Replaces TeX logo keywords with accessible HTML spans. \\
\texttt{texsmith.extensions.index} & Adds the \texttt{\#[tag]} syntax, LaTeX index entries and an MkDocs plugin. \\
\texttt{texsmith.extensions.counters} & Adds the \texttt{\#\{prefix:key\}} custom-counter syntax and an MkDocs plugin. \\

\bottomrule
\end{tabularx}
\end{center}

Inspect the pipeline's default extension list programmatically:

\begin{code}{python}{}{}
\begin{Verbatim}[commandchars=\\\{\},breaklines, breakanywhere, commandchars=\\\{\}]
\PY{o}{\PYZgt{}\PYZgt{}}\PY{o}{\PYZgt{}} \PY{k+kn}{from}\PY{+w}{ }\PY{n+nn}{texsmith}\PY{n+nn}{.}\PY{n+nn}{adapters}\PY{n+nn}{.}\PY{n+nn}{markdown}\PY{+w}{ }\PY{k+kn}{import} \PY{n}{DEFAULT\PYZus{}MARKDOWN\PYZus{}EXTENSIONS}
\PY{o}{\PYZgt{}\PYZgt{}}\PY{o}{\PYZgt{}} \PY{p}{[}\PY{n}{e} \PY{k}{for} \PY{n}{e} \PY{o+ow}{in} \PY{n}{DEFAULT\PYZus{}MARKDOWN\PYZus{}EXTENSIONS} \PY{k}{if} \PY{n}{e}\PY{o}{.}\PY{n}{startswith}\PY{p}{(}\PY{l+s+s2}{\PYZdq{}}\PY{l+s+s2}{texsmith.}\PY{l+s+s2}{\PYZdq{}}\PY{p}{)}\PY{p}{]}\PY{p}{[}\PY{p}{:}\PY{l+m+mi}{2}\PY{p}{]}
\PY{p}{[}\PY{l+s+s1}{\PYZsq{}}\PY{l+s+s1}{texsmith.extensions.index:TexsmithIndexExtension}\PY{l+s+s1}{\PYZsq{}}\PY{p}{,} \PY{l+s+s1}{\PYZsq{}}\PY{l+s+s1}{texsmith.extensions.multi\PYZus{}citations:MultiCitationExtension}\PY{l+s+s1}{\PYZsq{}}\PY{p}{]}
\end{Verbatim}
\end{code}
\section{Using the extensions with Python Markdown}\label{using-the-extensions-with-python-markdown}

Pass the \texttt{module:attribute} strings (the same form used in
\texttt{DEFAULT\_MARKDOWN\_EXTENSIONS}) straight to Python Markdown:

\begin{code}{python}{}{}
\begin{Verbatim}[commandchars=\\\{\},breaklines, breakanywhere, commandchars=\\\{\}]
\PY{k+kn}{from}\PY{+w}{ }\PY{n+nn}{markdown}\PY{+w}{ }\PY{k+kn}{import} \PY{n}{Markdown}

\PY{n}{md} \PY{o}{=} \PY{n}{Markdown}\PY{p}{(}
    \PY{n}{extensions}\PY{o}{=}\PY{p}{[}
        \PY{l+s+s2}{\PYZdq{}}\PY{l+s+s2}{texsmith.extensions.smallcaps:SmallCapsExtension}\PY{l+s+s2}{\PYZdq{}}\PY{p}{,}
        \PY{l+s+s2}{\PYZdq{}}\PY{l+s+s2}{texsmith.extensions.texlogos:TexLogosExtension}\PY{l+s+s2}{\PYZdq{}}\PY{p}{,}
        \PY{l+s+s2}{\PYZdq{}}\PY{l+s+s2}{texsmith.extensions.index:TexsmithIndexExtension}\PY{l+s+s2}{\PYZdq{}}\PY{p}{,}
    \PY{p}{]}
\PY{p}{)}
\PY{n}{html} \PY{o}{=} \PY{n}{md}\PY{o}{.}\PY{n}{convert}\PY{p}{(}\PY{l+s+s2}{\PYZdq{}}\PY{l+s+s2}{`\PYZsh{}[LaTeX]` renders a TeX logo and an index entry.}\PY{l+s+s2}{\PYZdq{}}\PY{p}{)}
\end{Verbatim}
\end{code}
If you prefer the explicit class names the following also works:

\begin{code}{python}{}{}
\begin{Verbatim}[commandchars=\\\{\},breaklines, breakanywhere, commandchars=\\\{\}]
\PY{k+kn}{from}\PY{+w}{ }\PY{n+nn}{texsmith}\PY{n+nn}{.}\PY{n+nn}{extensions}\PY{n+nn}{.}\PY{n+nn}{index}\PY{+w}{ }\PY{k+kn}{import} \PY{n}{TexsmithIndexExtension}

\PY{n}{md} \PY{o}{=} \PY{n}{Markdown}\PY{p}{(}\PY{n}{extensions}\PY{o}{=}\PY{p}{[}\PY{n}{TexsmithIndexExtension}\PY{p}{(}\PY{p}{)}\PY{p}{]}\PY{p}{)}
\end{Verbatim}
\end{code}
\section{Register extensions in MkDocs}\label{register-extensions-in-mkdocs}

Once TeXSmith is installed you can reference the modules directly from
\texttt{mkdocs.yml}:

\begin{code}{yaml}{}{}
\begin{Verbatim}[commandchars=\\\{\},breaklines, breakanywhere, commandchars=\\\{\}]
\PY{n+nt}{markdown\PYZus{}extensions}\PY{p}{:}
\PY{+w}{  }\PY{p+pIndicator}{\PYZhy{}}\PY{+w}{ }\PY{l+lScalar+lScalarPlain}{texsmith.index}
\PY{+w}{  }\PY{p+pIndicator}{\PYZhy{}}\PY{+w}{ }\PY{l+lScalar+lScalarPlain}{texsmith.texlogos}
\PY{+w}{  }\PY{p+pIndicator}{\PYZhy{}}\PY{+w}{ }\PY{l+lScalar+lScalarPlain}{texsmith.extensions.smallcaps}
\end{Verbatim}
\end{code}
(\texttt{texsmith.index} and \texttt{texsmith.texlogos} are registered entry-point aliases for
\texttt{texsmith.extensions.index} / \texttt{texsmith.extensions.texlogos}.)

The index extension also publishes an MkDocs plugin that injects collected tags
into the \texttt{search\_index.json}. Enable it next to the Markdown extension:

\begin{code}{yaml}{}{}
\begin{Verbatim}[commandchars=\\\{\},breaklines, breakanywhere, commandchars=\\\{\}]
\PY{n+nt}{plugins}\PY{p}{:}
\PY{+w}{  }\PY{p+pIndicator}{\PYZhy{}}\PY{+w}{ }\PY{l+lScalar+lScalarPlain}{texsmith.index}
\end{Verbatim}
\end{code}
When you use \texttt{mkdocs-\allowbreak{}texsmith} the plugin automatically appends both the
Markdown extension and the MkDocs plugin unless you disable it with the
\texttt{inject\_markdown\_extension} option.

\section{Integrating with the LaTeX renderer}\label{integrating-with-the-latex-renderer}

TeXSmith renders through a typed IR: \texttt{read(HTML) \(\rightarrow\) IR \(\rightarrow\) write(IR) \(\rightarrow\) LaTeX}.
Extensions that need \LaTeX{} output add a reader lowering (\texttt{@reads}, HTML \(\rightarrow\) IR)
and a writer emitter (\texttt{@writes}, IR \(\rightarrow\) \LaTeX{}) instead of mutating HTML. See
\hyperref[snippet:<HASH>]{Readers \& Writers} for the decorators and a complete,
runnable example (\texttt{examples/custom-\allowbreak{}render/counter.py}).

\section{Write your own extensions}\label{write-your-own-extensions}

You can create custom Markdown extensions that plug into TeXSmith's conversion
pipeline. Refer to the API documentation for details on
the extension points and how to register your extension with TeXSmith.

\subsection{Pipeline placement \& precedence}\label{pipeline-placement-precedence}

\begin{itemize}
\item{} Markdown extensions run before slot extraction and fragment rendering; any HTML they emit flows through the same pipeline.
\item{} Reader lowerings (\texttt{@reads}) are selected by \texttt{(level, tag)} and priority; the resulting IR is then emitted by writer emitters (\texttt{@writes}) dispatched on node type.
\item{} Extensions should not override template partials directly---expose fragment partials or attributes instead so precedence stays transparent.

\end{itemize}

\end{document}
